\section{Giới thiệu}

Trong giai đoạn chuyển đổi số, các khóa học trực tuyến đại trà đã trở thành công cụ thiết yếu để xây dựng một xã hội học tập suốt đời. Sự thành công của một nền tảng học trực tuyến vừa phụ thuộc vào chất lượng học liệu số, vừa chịu sự chi phối bởi hiệu suất của hệ thống truyền thông hỗ trợ học viên.

\subsection{Bối cảnh nền tảng daotao.ai và Trung tâm Công nghệ Giáo dục}

Nền tảng học trực tuyến daotao.ai được phát triển và vận hành bởi Trung tâm Công nghệ Giáo dục thuộc Trường Công nghệ Thông tin và Truyền thông – Đại học Bách khoa Hà Nội. Nền tảng ra đời với nhiệm vụ cung cấp các khóa học đại trà phục vụ cộng đồng, đồng thời chịu trách nhiệm đào tạo và bồi dưỡng kỹ năng số cho cán bộ công chức trên toàn quốc theo Đề án 06 của Chính phủ. Với quy mô người học lớn và có sự phân hóa sâu sắc về nhân khẩu học cũng như trình độ công nghệ, Trung tâm Công nghệ Giáo dục phải đối mặt với áp lực vận hành hệ thống hỗ trợ kỹ thuật lớn.

\subsection{Tầm quan trọng của truyền thông tích hợp trong đào tạo trực tuyến}

Truyền thông trong giáo dục số là quá trình thiết kế truyền thông có mục tiêu nhằm duy trì sự tương tác ba chiều giữa giảng viên, học viên và học liệu. Một hệ thống truyền thông tích hợp hiệu quả sẽ:
\begin{itemize}
    \item Bảo đảm sự phối hợp nhịp nhàng và chính xác trong khâu số hóa học liệu giữa bộ phận kỹ thuật của trung tâm và giảng viên chuyên môn tại Đại học Bách khoa Hà Nội.
    \item Giúp người học nhanh chóng tiếp cận hướng dẫn học tập, vượt qua các rào cản thao tác kỹ thuật ban đầu để tập trung vào việc tiếp thu tri thức.
    \item Cung cấp vòng phản hồi nhanh chóng nhằm nâng cao trải nghiệm học tập trực tuyến, củng cố lòng tin và giảm tỷ lệ học viên bỏ học giữa chừng.
\end{itemize}

\subsection{Các bất cập thực tế và mục tiêu nghiên cứu}

Mặc dù đã triển khai nhiều cải tiến, hệ thống truyền thông hiện tại của daotao.ai vẫn bộc lộ những bất cập:
\begin{itemize}
    \item \textbf{Trong nội bộ và đối tác:} Việc lệ thuộc vào nhóm chat Zalo tự phát để phối hợp thông số kỹ thuật học liệu dẫn đến hiện tượng trôi tin nhắn và hiểu sai thông điệp, gây lãng phí nguồn lực biên soạn bài giảng và học liệu chuẩn SCORM.
    \item \textbf{Trong tương tác với học viên:} Tình trạng học viên, đặc biệt là cán bộ Đề án 06 lớn tuổi, gặp khó khăn khi tìm nút đăng ký học, không tự khắc phục được sự cố đăng nhập thông thường kết hợp với tốc độ phản hồi thủ công chậm trễ của trung tâm đã tạo ra rào cản nhận thức cho người học.
\end{itemize}

Để khắc phục các tồn tại trên, báo cáo này thực hiện phân tích hệ thống truyền thông hiện tại của daotao.ai dựa trên bốn lý thuyết truyền thông gồm mô hình truyền thông tuyến tính Shannon--Weaver \parencite{shannon1949}, mô hình giao dịch \parencite{barnlund2008}, thuyết phong phú truyền thông \parencite{daft1986} và thuyết tải nhận thức \parencite{sweller1988}. Từ đó, báo cáo đề xuất và kiểm chứng thực nghiệm hai giải pháp cải tiến thực tế gồm tích hợp chatbot tự động trợ giúp trên website và thiết kế bộ mẫu email trực quan, hướng tới tối ưu hóa hiệu suất vận hành giáo dục số của Trung tâm Công nghệ Giáo dục.
